%% CPP 2027 Submission — 0pat Exercise XXV: Orthogonal Euler Circles as an Observation Frame
\documentclass[sigplan,10pt,anonymous,review]{acmart}
\settopmatter{printfolios=true,printacmref=false}

\AtBeginDocument{%
  \providecommand\BibTeX{{\normalfont B\scshape ib\TeX}}}

\settopmatter{printacmref=false}
\renewcommand\footnotetextcopyrightpermission[1]{}
\pagestyle{plain}

\usepackage{microtype}
\usepackage{booktabs}
\usepackage{amsmath}
% XeTeX (tectonic): acmart loads newtxmath/libertine which already defines \Bbbk;
% reset it so amssymb's definition is accepted.
\let\Bbbk\relax
\usepackage{amssymb}

\begin{document}

\title{Orthogonal Euler Circles as an Observation Frame:\\ The Chinese
Remainder Theorem and Prime Observability (Exercise XXV)}

\author{Anonymous}
\affiliation{%
  \institution{Anonymous}
}
\email{anonymous@example.org}

\begin{abstract}
We report the third completed exercise of a re-formalization program
(``0pat''): the insight that the prime axis can be observed through two
orthogonal Euler circles --- the periodized number line $\mathbb{Z}/p\mathbb{Z}$
as a circle, two such circles at coprime moduli as an orthogonal coordinate
system --- is re-stated in pure known mathematics and formally proved in
Lean 4 / mathlib with \emph{no} pat framework concepts and \emph{no}
complex plane. The re-statement is the algebra of double-modulus
observation: two orthogonal circles ($\mathbb{Z}/p\mathbb{Z} \times
\mathbb{Z}/q\mathbb{Z}$) compose into the circle of the product modulus
by the Chinese remainder theorem
(\texttt{ZMod.chineseRemainder}); a prime factor $p$ of $a$ is visible
as the zero projection on the $\mathbb{Z}/p\mathbb{Z}$ circle; distinct
primes are distinguished by nonzero projection on each other's circles;
a prime observed on a non-own circle is a unit of that circle's group,
satisfying Fermat's little theorem $p^{q-1} \equiv 1 \pmod q$; and on
its own circle a prime satisfies the complete Fermat algebra
$a^p \equiv a \pmod p$. Seven theorems, zero errors, zero
\texttt{sorry}, mathlib v4.32.2.
\end{abstract}

\keywords{Lean, formalization, Chinese remainder theorem, Fermat's little theorem, circle group, 0pat, reproducibility}

\maketitle

\section{The observation frame}

The starting insight of this exercise is that the prime axis can be
observed through circles. The number line, periodized modulo $p$, is the
circle $\mathbb{Z}/p\mathbb{Z}$; the mapping between the real axis and a
circle is already formalized in pure known mathematics (stereographic
projection $\mathbb{R} \cup \{\infty\} \cong S^1$ and the covering map
$\mathbb{R} \to \mathbb{R}/2\pi\mathbb{Z}$), so observing integers on a
circle is a legitimate coordinate choice, not a new structure. The
question posed here is whether \emph{two} such circles, orthogonal to
each other, form an observation frame for the prime axis.

The answer, in pure known mathematics, is the algebra of double-modulus
observation. Two circles at moduli $p$ and $q$ are orthogonal exactly
when $p$ and $q$ are coprime: the product circle
$\mathbb{Z}/p\mathbb{Z} \times \mathbb{Z}/q\mathbb{Z}$ then composes
into the single circle $\mathbb{Z}/pq\mathbb{Z}$ of the product modulus.
The composition is the Chinese remainder theorem, present in mathlib as
\texttt{ZMod.chineseRemainder}: for coprime $m, n$, the rings
$\mathbb{Z}/mn\mathbb{Z}$ and $\mathbb{Z}/m\mathbb{Z} \times
\mathbb{Z}/n\mathbb{Z}$ are isomorphic. This is the coordinate law of the
frame: a number $a$ has exactly one position in each circle, and the pair
of positions determines $a$ modulo the product.

The frame is legitimate because the composition law is an isomorphism:
observing through two orthogonal circles loses nothing that observing
through the product circle would keep. The frame is useful because prime
structure becomes visible in the projections: this is the content of the
theorems below.

\section{Prime observability in the frame}

Seven theorems, all machine-checked, state what the frame sees. Each
statement is classical mathematics; the novelty is the exact formulation
of the observation and the machine-checked proof.

\paragraph{The coordinate law (orthogonal\_frame\_crt).} For coprime
$m, n$, the product of the two orthogonal circles is the circle of the
product modulus: \texttt{ZMod.chineseRemainder} gives
$\mathbb{Z}/(mn)\mathbb{Z} \simeq \mathbb{Z}/m\mathbb{Z} \times
\mathbb{Z}/n\mathbb{Z}$ as rings. This is the law that makes the frame a
coordinate system: two orthogonal circle positions determine one position
on the product circle, and conversely.

\paragraph{Prime factors are visible (factor\_observable\_zero).} If $p$
is prime and $p \mid a$, then $a$ projects to $0$ on the
$\mathbb{Z}/p\mathbb{Z}$ circle: $(a : \mathbb{Z}/p\mathbb{Z}) = 0$. A
prime factor of $a$ is therefore directly observable as a zero in the
corresponding coordinate --- the first coordinate of the frame detects
the factor.

\paragraph{Distinct primes are distinguished (distinct\_primes\_observable).}
If $p$ and $q$ are distinct primes, then $p$ projects to a nonzero
element of the $\mathbb{Z}/q\mathbb{Z}$ circle. The proof is
divisibility-theoretic: $q \nmid p$ (a prime divides a prime only if
they are equal), hence $p \bmod q \ne 0$. The frame therefore assigns
each prime a position that no other prime occupies: prime positions in
the frame are unique.

\paragraph{Observation on a non-own circle (prime\_unit\_on\_other\_circle).}
A prime $p$ observed on a circle of a different prime modulus $q$ is a
unit of that circle's group: $p^{q-1} \equiv 1 \pmod q$, by Lagrange's
theorem applied to the multiplicative group
$(\mathbb{Z}/q\mathbb{Z})^\times$ of order $q-1$ (the circle-group size
theorem, reused from Exercise XXIV). The prime is a regular point of the
foreign circle --- it occupies a group element, not the zero --- and its
position is governed by the circle's group law.

\paragraph{Observation on its own circle (fermat\_little\_complete).}
On the circle of its own modulus, a prime satisfies the complete Fermat
algebra: for every $a$, $a^p \equiv a \pmod p$. This is the full form of
Fermat's little theorem (coprime case via the unit group and Lagrange;
non-coprime case via congruence lifting), reused from Exercise XXIV.
Together with the previous paragraph: on its own circle the prime is the
modulus itself (everything is governed by $p$); on every other circle it
is a unit subject to that circle's group.

\section{Relation to the exercise series}

This exercise connects the two previous ones. Exercise XXIV established
the prime circle: a prime $p$ is the modulus of the circle
$\mathbb{Z}/p\mathbb{Z}$ and its algebra is Fermat's little theorem.
Exercise XXIII established that composites are prime factorizations. The
orthogonal frame completes the picture: a composite $n = pq$ is observed
as the product circle $\mathbb{Z}/pq\mathbb{Z}$, which the coordinate law
factors into the two orthogonal circles $\mathbb{Z}/p\mathbb{Z} \times
\mathbb{Z}/q\mathbb{Z}$; the prime factors are the zero-projections of
the frame (factor observability), and their circle-algebra is Fermat's
theorem. The frame is the geometric form of the Chinese remainder
theorem, and the observability theorems are its prime-theoretic reading.

The honest boundary of the exercise is stated as in the previous ones:
every theorem is a restatement of classical mathematics
(\texttt{novelty\_status: KNOWN}), the proofs are machine-checked with
zero \texttt{sorry} against mathlib v4.32.2, and no claim is made beyond
the statements listed. In particular, the frame changes coordinates; it
does not change the distribution of primes.

\section{Formalization}

The proof is a single Lean file (\texttt{proof.lean}, 272 lines),
importing \texttt{Mathlib.Data.ZMod.Basic},
\texttt{Mathlib.GroupTheory.OrderOfElement}, and
\texttt{Mathlib.Data.Nat.Prime.Basic}. The seven theorems:

\begin{center}
\begin{tabular}{lll}
\toprule
Theorem & Statement & mathlib anchor \\
\midrule
\texttt{orthogonal\_frame\_crt} & $m \perp n \Rightarrow \mathbb{Z}/mn \simeq \mathbb{Z}/m \times \mathbb{Z}/n$
  & \texttt{ZMod.chineseRemainder} \\
\texttt{factor\_observable\_zero} & $p \mid a \Rightarrow (a : \mathbb{Z}/p) = 0$
  & \texttt{ZMod.val\_injective}, \texttt{mod\_eq\_zero\_of\_dvd} \\
\texttt{distinct\_primes\_observable} & $p \ne q \Rightarrow (p : \mathbb{Z}/q) \ne 0$
  & \texttt{Nat.Prime.eq\_one\_or\_self\_of\_dvd} \\
\texttt{card\_units\_zmod\_prime} & $|(\mathbb{Z}/p)^\times| = p-1$
  & explicit bijection (Ex. XXIV) \\
\texttt{fermat\_little} & $\gcd(a,p)=1 \Rightarrow a^{p-1} \equiv 1$
  & \texttt{orderOf\_dvd\_card} \\
\texttt{fermat\_little\_complete} & $\forall a,\ a^p \equiv a \pmod p$
  & \texttt{Nat.ModEq.pow} \\
\texttt{prime\_unit\_on\_other\_circle} & $p \ne q \Rightarrow p^{q-1} \equiv 1 \pmod q$
  & Lagrange + CRT-free unit group \\
\bottomrule
\end{tabular}
\end{center}

Compiled with Lean 4.32.2 / mathlib v4.32.2, zero errors, zero
\texttt{sorry}, zero new axioms. The file is self-contained; the
circle-group-size and Fermat theorems are restated inside it (reused
from Exercise XXIV) so that the exercise does not depend on any other
exercise file.

\section{Conclusion}

The orthogonal Euler frame is a legitimate coordinate system for
observing the prime axis: its coordinate law is the Chinese remainder
theorem, its prime observations are Fermat's little theorem and the
zero/nonzero projection theorems, and all of it is machine-checked in
pure known mathematics. The frame is the geometric form of the CRT, and
it connects the composite factorization of Exercise XXIII with the prime
circle of Exercise XXIV. Coordinates change; distributions do not ---
the honest boundary is kept, and the frame's value is observational
clarity, not new number theory.

\bibliographystyle{ACM-Reference-Format}
\bibliography{refs}

\end{document}
